\documentclass[a4paper,10pt]{article}
\newcommand\subdir{/home/koen/LaTeX-setup/}
\input{\subdir/LaTeX-files/imports.tex}
\input{\subdir/LaTeX-files/master-internship.tex}
\usepackage{\subdir/LaTeX-files/aastex}
\usepackage{natbib}
\bibliographystyle{\subdir/LaTeX-files/astroadstitle.bst}

%Set the title, Author and date
\title{Notes Week 28}
\author{Koen Essers}
\tutorial{Master Internship}
\date{\today}

%Set the header style
\header{\tutorialstring}

%Begin the document
\begin{document}
\setuplayout
\section{Monash Initial Helium Envelope Abundances}

Last week, I showed these two plots:

\begin{figure}[ht]
    \marginfignote{0}{The most important thing we saw here was the quite large differences between the helium mass fractions.}
    \centering
    \incplot{w27-initial-envelope-abundance}
\end{figure}

\begin{figure}[ht]
    \marginfignote{0}{Here, the same difference is again very clear.}
    \centering
    \incplot{w27-compare-he}
\end{figure}

It is actually very easy to understand where this discrepancy comes from, and in fact, I have already made a figure before that shows exactly what is going wrong here:

\begin{figure}[ht]
    \centering
    \incplot{w26-initial-envelope-zoom}
\end{figure}
Here I plot the initial envelope abundances for 3 different Monash metallicities divided by the Monash solar Metallicity model.

I made this plot to \emph{confirm} I can simply take the abundance of an element, divide it by the Monash metallicity and multiply it by the MESA metallicity to get the MESA envelope abundance.

Now I think this works great for pretty much all elements, \emph{except} for some of the very light elements, exactly the elements that I was comparing before.

I quickly fixed this by also interpolating the initial envelope abundances.

\begin{figure}[ht]
    \marginfignote{5}{Clearly, this works much better for helium and actually all other isotopes too.}
    \centering
    \incplot{w28-initial-envelope-abundance}
\end{figure}

\begin{figure}[ht]
    \marginfignote{0}{This looks much better!}
    \centering
    \incplot{w28-compare-he}
\end{figure}



\newsection{MESA intershell abundances}

In order to verify the Monash enrichment is working well, we can compare the enrichment of the envelope as predicted by my Monash post-processing code with the abundances from MESA.

I think there are a couple interesting things to do. First, we just want to compare the intershell abundances of the elements that are available to me in the MESA models with the
predicted Monash intershell abundances. This is already interesting, because this gives us a good idea of it the interpolation is working well.

Then, I think it is interesting to look at the envelope abundances to compare the differences. I have already done this previously, \emph{but} now I have a ``better'' initial envelope abundance, and insight into the intershell abundances of both codes.

Lastly, I want to verify that my envelope enrichment code is working correctly. I will do this by injecting the MESA intershell abundances and the initial envelope abundance into
my post-processor. I will then compare the envelope abundances it spits out with the actual MESA envelope abundances. If the post-processor is working well, it should reproduce the
MESA envelope abundances well\sidenote{Although I do not expect it to agree completely, because the dredge up assumed to be instantaneous, and the precise derived dredged up masses
are slightly different from the MESA internal dredged-up masses.}.

\subsection{Intershell abundances}

To start, I document \emph{how} I obtain the intershell abundances using only a few profiles. 
\begin{figure}[ht]
    \marginfignote{0}{For all single stars, I save a model and profile at the start of a thermal pulse, just before things really start to change and the helium burning is just beginning to start up.}
    \marginfignote{8}{These profiles contain a lot of relevant information, but they are sparse. To see \emph{when} in the lifetime of stars the profiles are saved, I show the core mass and
    the helium luminosity.}
    \marginfignote{16}{Clearly, the models are correctly saved at the start of a thermal pulse. I think this is quite a nice time to sample the intershell abundance, since this is roughly also the abundance when dredge up occurs.}
    \marginfignote{26}{The last profile is an outlier. This is the model that is saved when the envelope mass is below \(0.01\; M_\odot \).}
    \centering
    \incplot{w28-profiles-1}
\end{figure}

Now let me plot something a bit more interesting. By finding the position of the greatest change in the mass fraction of the metals\sidenote{over a small number of zones},
\begin{minted}[fontsize=\small,breaklines]{py}
i = np.argmax(
    np.abs(
        np.log10(p.z_mass_fraction_metals[10:])
        - np.log10(p.z_mass_fraction_metals[:-10])
    )
)
\end{minted}
we get a proxy for the intershell / envelope boundary.

By then simply taking the mass averaged abundance of an element a certain margin before this boundary, we get the envelope abundance.
Since we also have my MESA implementation of the envelope abundances, we can compare the two.

\begin{figure}[ht]
    \marginfignote{0}{We clearly see that the two agree very well, even though they are computed with slightly different methods.}
    \marginfignote{5}{The MESA history data is based on the largest convective region. It then takes a mass average of this region. Luckily, and also unsurprisingly, this agrees well with the method that simply takes the mass averaged value of the abundances past the intershell zone.}
    \marginfignote{16}{The only difference occurs when the envelope mass is very small.}
    \centering
    \incplot{w28-profiles-envelope-check}
\end{figure}

The shape of the envelope abundance, and very likely also the intershell abundance, is very nicely approximated by a step function.

Using the same criterion, but now taking some average to the intershell side of the intershell / envelope boundary, we can also find the intershell abundances.
Unfortunately however, we cannot compare this with \emph{MESA} history data.

\begin{figure}[ht]
    \marginfignote{0}{Here we can see all available elements together to see the relative abundances.}
    \marginfignote{6}{Everything looks normal to me, except for weird nitrogen bump. The nitrogen abundance is constantly }
    \centering
    \incplot{w28-profiles-intershell}
\end{figure}

\begin{figure}[ht]
    \marginfignote{0}{This is weird I think. But it is not necessarily something that is wrong with my
    intershell abundance finding code. There is just something that happens during this specific thermal pulse
that makes the intershell relatively \emph{very} abundant.}
    \centering
    \incplot{w28-profiles-nitrogen}
\end{figure}

\begin{figure}[ht]
    \marginfignote{0}{Apparently, the \(1.5\; M_\odot \) model just by accident, only has 1 thermal pulse with a significant nitrogen abundance.
    This however does not mean that all stars have this 1 thermal pulse with a large nitrogen abundance, and it also does not seem like an error since there is quite a clear pattern.}
    \marginfignote{14}{Clearly, stars with masses larger than \(1.5\; M_\odot \) have multiple pulses with a large
    nitrogen abundance. Stars with masses less than \(1.5\; M_\odot \) have \emph{no} thermal pulses with a
large nitrogen abundance.}
    \marginfignote{23}{I do not know what causes these sudden increases in envelope abundances and why the
    abundances drop to low values again during the last thermal pulse.}
    \marginfignote{31}{Maybe it has something to do with the CNO cycle?}
    \centering
    \incplot{w28-profiles-intershell-nitrogen}
\end{figure}

\nextpage
Obviously, the mass is quite important for the nitrogen abundance. Let me also plot the intershell
abundances for the \(3\; M_\odot \) MESA model.

\begin{figure}[ht]
    \marginfignote{0}{Obviously the nitrogen abundance is quite different for this \(3\; M_\odot \) model, but
    we see differences in the other isotopes too. Especially the neon abundance is much larger in this model.}
    \marginfignote{10}{Carbon seems \emph{less} abundant in the heavier model.}
    \centering
    \incplot{w28-profiles-intershell-m3}
\end{figure}

This figure below compares the MESA intershell abundances with the Monash post-processed
intershell abundances for the \(2.8\; M_\odot \) model.
\begin{figure}[ht]
    \centering
    \incplot{w28-profiles-intershell-m2.5-monash}
\end{figure}



\subsection{Method}
\subsection{Results}

\newsection{Investigating Monash Yields}
% why are the mesa abundances structurally about a factor 2 larger? and why is there a difference between elements before say element pr

\newsection{Envelope abundances}
% compare the envelope abundances per pulse of the Monash models from the paper with the MESA abundances.

\bibliography{master.bib}
\end{document}
